\section{Original quartic Type II exclusion in two weight progressions}
\label{sec:cubic-original-quartic-typeii-exclusion}

Put
\[
 \begin{gathered}
 p=269,\qquad K=\mathbb Q(\sqrt{-3}),\qquad
 \eta=\chi_{-3},\qquad N_\psi=10,\\
 \psi_5(2)=i,\quad\operatorname{cond}(\psi_5)=5,\qquad
 \xi=\eta\psi_5^2=\chi_{-15}.
 \end{gathered}
 \tag{G128.1}
\]
Here \(\psi\) is the character modulo \(10\) induced by the
primitive character \(\psi_5\); evaluation at \(2\) in (G128.1)
is in the primitive modulus \(5\).  Write
\(\Psi=\operatorname{Chr.dirG}(\psi)=\operatorname{Chr.dirG}(\psi_5)\),
so that \(\Psi(\operatorname{Frob}^{\rm ar}_r)=\psi_5(r)\) at
unramified primes.  We also use \(\xi\) for its quadratic Galois
character when it occurs in a representation formula.
The character \(c_p\) is the actual action on
\(p\)-power roots of unity.  All Galois Frobenius polynomials below
use geometric Frobenius.  Assume
\[
                  e\ge1,\qquad
                  e\equiv133\ \hbox{or}\ 267\pmod{268},
                  \qquad k=2e+1.
 \tag{G128.2}
\]
Thus \(e\) and \(k\) are odd, \(k\equiv267\pmod{268}\), and
\(268\nmid2e\).  In particular these residues respect the parity
and genericity clauses of the literal original datum.

\begin{theorem}[AC-G128-quartic-exclusion: the original quartic Type II set is empty]
\label{AC-G128-quartic-exclusion}
Retain the complete original datum \(\operatorname{Dat}\) for
\((K,e,N_\psi,\psi)\), the actual cubic ring and selected local
prime of AC-G65, and every hypothesis used by AC-G68-point.
In particular retain \(\mathrm{GD}\), \(\mathrm{CohInp}\),
\texttt{AutomorphyUnramifiedCoefficientsAlt}, the original
coefficient condition \(e_{O_0}=1\), and the published local
compatibility and primitive-newform inputs used in AC-G116.
For either \(j=1,2\), there is no actual characteristic-zero Type II
point of this original cubic model at any weight in (G128.2).

The contradiction applies more generally to any actual integral
irreducible odd representation \(V:G_{\mathbb Q}\to
\operatorname{GL}_2(E)\), for a finite value field \(E\), having
the following precise properties supplied by AC-G68:
\[
 \begin{gathered}
 \det V=\eta,\qquad R_f=Vc_p^{-e},\\
 \overline V^{\rm ss}
   =(\eta\overline\Psi\,\overline c_p^{\,e})
       \oplus(\overline\Psi^{-1}\overline c_p^{-e}).
 \end{gathered}
 \tag{G128.3}
\]
where \(R_f\) is attached to an actual cuspidal newform \(f\) of weight \(k\);
the original exact away-inertia condition holds; and the actual
ordered ordinary, selected potentially crystalline condition at
\(p\) holds.  The statement therefore also excludes a native old
Type II endpoint with these properties directly.  No reverse
embedding of an arbitrary old point into the cubic selected
component is required.
\end{theorem}

\begin{proof}
\proofstep{1. Obtain the exact endpoint, with the original coefficient datum.}
Suppose first that an actual cubic Type II point exists.
AC-G68-block makes its middle character one, its endpoint
determinant \(\chi^{2j}\), and its endpoint block
\(\chi^jV|_{G_K}\).  AC-G68-point removes \(\chi^j\), proves
membership in the original selected ideal \(I_t(D)\), checks all
seven fields of \(D.\mathrm{IsPt}\), and obtains the actual old
modular endpoint \(V\).  Its proof uses \(O_0+\mathfrak m_E\)
when a finite value-field extension is needed.  It leaves
\(O_0\), its residue field, and the retained condition \(e_{O_0}=1\)
unchanged.  Its output is (G128.3), not merely a formal residual
trace with unspecified local conditions.

For a native old endpoint one instead uses
\texttt{typeII\_facts} and its explicit \texttt{Wfun} formula in
\texttt{CoboundaryLemmaNoColumn.lean}, followed by the full
\texttt{admissible\_of\_typeII} or its dense modular-extension
variant.  These give the same representation-level conditions.
The remainder of the proof uses only those conditions.
The original source \texttt{Level} parameter is \(10|D_K|=30\);
we do not identify that parameter with a primitive conductor.

\proofstep{2. Exact inertia, crystallinity, and the primitive level-fifteen twist.}
The full source \texttt{pt\_inertia\_away} proof, or the
prime-to-\(p\) residual-kernel projectors in AC-G68-point, gives
the actual inertia characters \(\Psi,\Psi^{-1}\) on \(G_K\)
away from \(p\).  The action is finite, so its monodromy is zero.
At every rational \(r\ne3,p\), all rational inertia is in
\(G_K\), since \(K/\mathbb Q\) is unramified at \(r\).
Consequently \(V|_{I_5}\) has characters
\(\Psi,\Psi^{-1}\), both of conductor exponent one.  At
\(r\ne3,5,p\) it is unramified, including at \(r=2\).
The factor \(2\) in the displayed modulus \(N_\psi=10\)
does not create inertia for the primitive conductor-five character.

At \(3\), \(\Psi\) is unramified, as follows also from
\(\mathrm{GD.cop}\).  Inertia over \(K_3\) acts trivially.
Rational inertia therefore acts through \(C_2\); its determinant
\(\eta\) forces the two characters \(1,\eta\).  This is finite
tame inertia, again with zero monodromy.

At \(p\), the literal character in \(D.\mathrm{loc}\) is
\((\operatorname{Chr.dirG}(\operatorname{comp}_p\psi))\circ D.\iota\).
It is one here because \(\psi\) has conductor \(5\).
AC-G68-point gives crystallinity over the unramified quadratic
\(K_p/\mathbb Q_p\) and trivial potentially crystalline inertial
type.  Semilinear descent of the full crystalline period module
therefore makes \(V\) crystalline over \(\mathbb Q_p\).
Equivalently this is the unramified descent step in the audited
good-twist argument.  The complete native admissibility proof
constructs an ordered line on the rational decomposition group,
with inertia \(c_p^e\), and quotient inertia \(c_p^{-e}\).
Only the finite unramified factors are left unspecified.

The determinant of \(R_f\) makes its classical nebentypus
\(\eta\), by the good-prime attachment and character uniqueness.
Let \(g\) be the normalized \emph{primitive newform} of the
actual automorphic twist
\[
                 \pi_g=\pi_f\otimes((\eta\psi_5)\circ\det).
 \tag{G128.4}
\]
The convention \texttt{Newform.Att} evaluates the attached
representation at inverse arithmetic Frobenius.  Thus twisting
the classical eigenvalues by \(\eta(r)\psi_5(r)\) corresponds to
\[
 \begin{gathered}
       R_g=V\eta\Psi^{-1}c_p^{-e},\qquad
       \det R_g=\xi c_p^{-2e},\\
       \overline R_g^{\rm ss}=1\oplus\xi\overline c_p^{-2e}.
 \end{gathered}
 \tag{G128.5}
\]
The determinant identity uses
\(\Psi^{-2}=\operatorname{DirG}(\psi_5^2)=\operatorname{DirG}(\chi_5)\).
The identity for \(R_g\) follows from good-prime polynomials
and Chebotarev, using the retained classical attachment.
It is important here that the Galois twist is \(\Psi^{-1}\);
the equality \(\psi=\psi^{-1}\) used in AC-G117 is unavailable.
The automorphic twist preserves weight and cuspidality, and
its nebentypus is \(\eta\psi_5^2=\xi\).

At \(5\), the two inertia characters of \(R_g\) are
\(1,\Psi^{-2}= \chi_5\).  At \(3\) they are \(1,\eta\).
At every other finite \(r\ne p\) it is unramified.
At \(p\) it remains crystalline, since the finite twist is
unramified and the Tate twist is crystalline.
The full local--global compatibility inputs of AC-G116
identify these Frobenius-semisimple Weil--Deligne parameters
with the local automorphic factors.  The Artin conductor is
one at each of \(3,5\), zero elsewhere, and monodromy is zero.
The coefficient-prime local parameter is unramified and hence
spherical.  Classical primitive-newform conductor theory gives
\[
                         g\in S_k(\Gamma_0(15),\xi)
                 \quad\hbox{of exact primitive level }15.
 \tag{G128.6}
\]
For comparison the primitive conductor of the original \(f\)
is \(3\cdot5^2=75\): both quartic characters contribute one
at \(5\).  Thus we distinguish the displayed
\(N_\psi=10\), the formal source level \(30\), and the actual
primitive levels \(75\) and \(15\) before and after (G128.4).
No level-lowering congruence theorem and no coefficientwise
imprimitive Dirichlet twist has been used.

\proofstep{3. The actual roots at both bad primes and at the coefficient prime.}
For every prime \(r\nmid15p\), (G128.5) gives
\[
                    a_r(g)\equiv1+\xi(r)r^{k-1}
                                      \pmod{\mathfrak m_E}.
 \tag{G128.7}
\]
The missing primes cannot be recovered from this good-prime
identity alone.  We determine their actual primitive roots.

For \(r=3\) and \(r=5\), the image of inertia in \(R_g\) is
exactly \(C_2\), acting with characters \(1,\nu_r\), where
\(\nu_3=\eta|_{I_3}\) and \(\nu_5=\chi_5|_{I_5}\).
Choose \(s_r\) in inertia with nontrivial image.  On the stable
integral lattice inherited from (G128.3),
\[
                         e_r^+=\frac{1+R_g(s_r)}2
 \tag{G128.8}
\]
is an integral rank-one projector onto inertia invariants.
It commutes with the decomposition action because inertia is
normal and \(\operatorname{Aut}(C_2)=1\).
Reduction commutes with this projector.  In (G128.5) the
unique inertia-trivial constituent at \emph{either} prime is
the global character \(1\): \(\xi\) is ramified at both
primes and the cyclotomic character is unramified there.
Thus the actual invariant line's character reduces to \(1\)
on the full decomposition group, even if the global residual
lattice was not initially split.

In the cohomological geometric-Frobenius convention, the
primitive Euler factors are
\[
 \det(1-X\operatorname{Fr}_r\mid R_g^{I_r})
                        =1-a_r(g)X\qquad(r=3,5).
 \tag{G128.9}
\]
It follows that \(a_3(g)\equiv a_5(g)\equiv1\).
These are coefficients of the primitive newform in (G128.4);
coefficientwise twisting at \(3\) or \(5\) would instead erase
them and would be the wrong operation.

The ordered line of \(R_g|_{G_{\mathbb Q_p}}\) is unramified.
Its saturated intersection with the stable lattice reduces to
a nonzero line.  The other residual inertia character is
\(\overline c_p^{-2e}\ne1\), so this line reduces to the
constituent \(1\) on the full decomposition group in (G128.5).
Its crystalline Frobenius eigenvalue \(\alpha\) is an integral
unit with \(\alpha\equiv1\).  The other root of the actual
Hecke polynomial
\[
                    X^2-a_p(g)X+\xi(p)p^{k-1}
 \tag{G128.10}
\]
is \(\xi(p)p^{k-1}/\alpha\), of positive valuation.
Consequently \(a_p(g)\equiv1\), and the actual form is ordinary.
At \(p\) this argument uses crystalline Frobenius, not an
unramified Galois Frobenius for the full representation.

\proofstep{4. Match every positive Fourier coefficient.}
Use the normalized Eisenstein series
\[
 E_k(1,\xi)=c_k+\sum_{n\ge1}b_nq^n,\qquad
 c_k=\frac{L(1-k,\xi)}2,\qquad
 b_n=\sum_{d\mid n}\xi(d)d^{k-1}.
 \tag{G128.11}
\]
At \(r=3,5\), \(b_{r^m}=1\).
For \(r\nmid15\), both the primitive newform and the
Eisenstein coefficients obey the same prime-power recurrence
with determinant term \(\xi(r)r^{k-1}\), and both are
multiplicative.  The values (G128.7)--(G128.10), including
\(p\), therefore give
\[
                         a_n(g)\equiv b_n\pmod{\mathfrak m_E}
                              \qquad(n\ge1).
 \tag{G128.12}
\]
At \(p\) the determinant term vanishes after reduction; at
\(3,5\) the primitive recurrence is \(a_{r^m}=a_r^m\).
All coefficients used here are integral: this follows from
the integral good-prime traces, the invariant-line bad roots,
the two integral roots at \(p\), and the recurrences.

\proofstep{5. The all-weight constant is a unit.}
The exact arithmetic in the retained quartic weight-fork
diagnostic gives
\[
                  c_{267}\equiv c_{535}\equiv145\pmod{269}.
 \tag{G128.13}
\]
It was computed over the rationals by the defining Bernoulli
recurrence and independently by the Akiyama--Tanigawa
algorithm.  A short finite-field verification of the reference
value is also possible with
\[
 \begin{aligned}
 B_0&=1,\qquad
 B_m=-\frac1{m+1}\sum_{r<m}\binom{m+1}{r}B_r
                                  \quad(1\le m\le267),\\
 B_n(X)&=\sum_{r=0}^n\binom nr B_rX^{n-r},\\
 c_{267}&=-\frac{15^{266}}{2\cdot267}
                    \sum_{a=1}^{15}\xi(a)B_{267}(a/15)\\
         &=145\quad\hbox{in }\mathbb F_{269}.
 \end{aligned}
 \tag{G128.14}
\]
Every denominator in this calculation is prime to \(269\).
The new verification script records the character table,
all intermediate ordinary Bernoulli residues, and the result;
the two earlier exact-rational outputs are preserved.

To pass to every weight in (G128.2), use Katz's actual integral
measure \(2H^{2,2}\) on
\(\mathbb Z_p\times\mathbb Z/15\mathbb Z\).  Its moment
normalization and Theorem 3.3.3 give
\[
 H^{2,2}(x^{k-1}\xi(t))=(1-2^k\xi(2))E_k(1,\xi)
                                      =(1-2^k)E_k(1,\xi).
 \tag{G128.15}
\]
Here \(\xi(2)=1\).  Since \(2^{134}=-1\) in \(\mathbb F_{269}\),
an odd power of \(2\) is neither \(1\) nor \(-1\).
In particular the smoothing factor is a unit, and for these
weights it is exactly \(135\) modulo \(269\).
For \(k\equiv267\pmod{268}\), the functions
\(x^{k-1}\xi(t)\) and \(x^{266}\xi(t)\) are congruent on
the whole measure domain: on units use Fermat, and on
\(p\mathbb Z_p\) both powers reduce to zero.
Taking constants in the integral measure and dividing by
the equal smoothing units proves
\[
                              c_k\equiv145\pmod{269}.
 \tag{G128.16}
\]
This includes actual weights divisible by \(269\); no division
by the actual weight is made modulo \(p\).
Also \(\xi(p)=-1\), so ordinary stabilization would multiply
the constant by \(1+p^{k-1}\), another unit.  Stabilization
does not remove (G128.16) and is not needed for the argument.

\proofstep{6. Katz's constant-expansion contradiction.}
Pull the two level-fifteen forms back to full tame level \(30\).
Enlarge only the finite value field to contain the required
roots of unity.  The original coefficient ring and its
\(e_{O_0}=1\) condition are unchanged.
Katz's generalized modular-form ring on the chosen determinant
component admits arbitrary coefficient base change and an
injective expansion at a chosen cusp, as stated and proved
in the applied framework of AC-G116.
Equation (G128.15) makes the Eisenstein series integral in
that ring.  The integral expansion of \(g\) places it there
as well by the divided-congruence inclusion: clear a fixed
denominator, then use expansion injectivity modulo successive
powers of the coefficient uniformizer to divide back.

After reduction, (G128.12) and (G128.16) say that \(E_k-g\)
has the constant expansion \(145\), with every positive
coefficient zero.  Injection identifies it with that nonzero
constant generalized form.  Both original forms have the
same weight character \(a\mapsto a^k\); a constant generalized
form has weight zero.  Applying \(a=2\) yields
\[
                           2^k\cdot145=145
                               \quad\hbox{in the residue field}.
 \tag{G128.17}
\]
This is impossible because \(2^k\ne1\) and \(145\ne0\).
Equivalently this is the constant-expansion case of Katz's
Corollary 4.4.1, with no small-weight restriction.
The contradiction proves the theorem.
\end{proof}

\begin{corollary}[The conditional quartic M4 consequence]
\label{cor:G128-quartic-M4}
For the actual original cubic models in (G128.1)--(G128.2),
retain their structural hypotheses in AC-G65--AC-G66 and
AC-G113.  Assume, in addition, actual \(\mathrm{M2}\) and
an actual host and native modular support satisfying exactly
(G110.2), with every classicality and published adjoint
hypothesis of AC-G110.  Then every two-sided
characteristic-zero local point is Type I and rigid, and the
actual model satisfies the literal \(\mathrm{M4}\) condition.
No mixed Type II Selmer-vanishing or
\texttt{RigiditySecondCase} hypothesis is needed for this
conditional conclusion.
\end{corollary}
\begin{proof}
The actual normalized shape calculation in AC-G113 has the
four alternatives Type I, upper, lower and Type II.  At a
two-sided point the model signs exclude upper and lower.
The theorem just proved excludes Type II under precisely
the original quartic local and modular hypotheses.
Thus each such point is Type I.  The actual support (G110.2)
allows AC-G110-rigidity to kill its original rational adjoint
Bloch--Kato group, and AC-G66-BK puts every actual dual
cocycle in that group.  Each dual cocycle is a coboundary.

Here are the remaining derivation quantifiers, to avoid
invoking the quadratic-only statement of AC-G122.
Let \(I\) be the actual homogeneous equation ideal of the
graded model \(A\), let \(\zeta:A\to\overline{\mathbb Q}_p\)
be a two-sided point killing \(I\), and let \(D,D'\) be two
point derivations at \(\zeta\) annihilating \(I\).
Actual M2 and the model point property give a torus scalar
\(t\ne0\) and an actual local point \(f:S_j\to
\overline{\mathbb Q}_p\) with
\(\operatorname{act}(t)\zeta=fq_A\).
On degree \(d\), define \(E(a)=t^dD(a)\) and
\(E'(a)=t^dD'(a)\), extending by the direct-sum grading.
The product rule proves that these are derivations at
\(fq_A\); the homogeneous generators and the product rule
prove that they annihilate \(I\).  The actual model's
dual-point extension property gives dual maps on \(S_j\).
Both signs are preserved because every \(t^d\) is nonzero.
Their cocycles are therefore coboundaries by the preceding
Type I argument.

The exact \texttt{normCoc\_cob\_iff} calculation uses that
\(2\) and the marked diagonal difference are units and that
the actual representation is normalized.  It puts both
cobounding matrices in the single line
\(\overline{\mathbb Q}_p\operatorname{diag}(1,0,-1)\).
Consequently the two cocycles are linearly dependent.
Actual dual-entry determination makes the same combination
of dual derivatives vanish on \(S_j\), and hence on \(A\).
Cancelling \(t^d\) on each homogeneous piece gives
\(aD+a'D'=0\) with \(a,a'\) not both zero.
This is the complete definition of M4, as in the generic
source derivation argument.  It does not use a quadratic
character identity.
\end{proof}

\paragraph{What this exclusion does and does not supply.}
The mixed tangent obligation in AC-G68 is vacuous for the
actual original Type II set just excluded.  This is not a
vanishing theorem for \(H_f^1(K,V|_{G_K}\otimes\chi^j)\)
for arbitrary representations.  Actual M2 and (G110.2)
in the corollary are additional premises, not consequences
of the Bernoulli calculation or of an isolated cusp form.
No new quartic class-sector bound, support host, selected
curve containment, count, or H04.4 conclusion is asserted.

There is no conflict with AC-G118: its original quartic
seed has \(e=23\), \(k=47\), where the constant has valuation
one rather than being a unit.  Nor is there a conflict with
AC-G127 at the present weights.  Its auxiliary seed has
nontrivial order-\(269\) type at \(q=3229\); after the same
primitive twist its conductor contains \(q\), and its
residual-trivial ramification violates the exact original
away-kernel premise.  The level-fifteen reduction in
(G128.6) therefore does not apply to it.
The earlier weight-fork diagnostic is preserved unchanged;
it only blocked one cusp-lifting construction.  This
exclusion is a separate proof with the missing actual
endpoint, local, bad-prime and modular hypotheses checked.

\paragraph{Source and computational scope.}
For this argument the complete AC-G116 and AC-G117 proofs,
the full AC-G68 interface and original-point transfer were
reread.  Fresh complete native reads cover
\texttt{Counts.lean}, \texttt{RhoKApply.lean},
\texttt{PtInertiaAway.lean}, \texttt{PtIneL.lean},
\texttt{CoboundaryLemmaNoColumn.lean},
\texttt{AdmissibleOfTypeII.lean}, its dense variant, and
\texttt{Newform.lean}.  For M4 the literal definitions and
complete applied \texttt{OrbitVanishingLocal.lean},
\texttt{GradedModelDerivationConditionOfPtHyp.lean} and
\texttt{NormCob.lean} proofs were read.  No Lean execution
or new proof of the old modularity declaration is claimed.

The local comparison and primitive-newform inputs retain
the exact Carayol, Saito and Casselman primary scope in
AC-G116.  Saito's geometric-Frobenius Euler-factor and
full compatibility statements were freshly reread in
\S2a and \S2c, Theorem 3, of
\url{https://www.ms.u-tokyo.ac.jp/~t-saito/pp/suams.pdf}.
Their deep proofs are imported.
Katz's \S1.2 arbitrary-base-change and chosen-cusp
injection, its complete applied Key Lemmas, and
\S3.1--3.3 including the complete printed Theorem 3.3.3
proof were freshly read at
\url{https://web.math.princeton.edu/~nmk/old/eismeasure.pdf}.
The exact earlier full Corollary 4.4.1 audit is retained.
The manifest distinguishes the two prior exact-rational
constant computations from the new finite-field replay
and preserves every input and output with its hash.
